\section{HỌC MÁY (MACHINE LEARNING)}

\subsection{Khái niệm}
Là tập con của trí tuệ nhân tạo. Một lĩnh vực nhỏ của khoa học máy tính, nghiên cứu các thuật toán và mô hình thống kê cho phép máy tính thực hiện các nhiệm vụ mà không cần lập trình rõ ràng.

\subsection{Các cơ chế học}
\begin{enumerate}
    \item Supervised Learning: Dữ liệu huấn luyện có nhãn (Input $X$, Output $Y$). Ví dụ: Phân loại email spam.
    \item Unsupervised Learning: Dữ liệu không có nhãn, máy tự tìm cấu trúc. Ví dụ: Phân cụm khách hàng.
    \item Reinforcement Learning: Học thông qua thử-sai và nhận thưởng/phạt.
\end{enumerate}

\subsection{Liên hệ với Khai phá dữ liệu}
{\centering\fontsize{17}{21}\selectfont\bfseries\boldmath
$\text{Data Mining} = \text{ML} + \text{Database} + \text{Statistics}$. \par}
Học máy cung cấp các công cụ và thuật toán kỹ thuật để thực hiện các bước Khai phá dữ liệu. Data Mining chú trọng vào quy trình tìm tri thức, Machine Learning chú trọng vào thuật toán dự đoán/học.

\subsection{Ứng dụng của Học máy}
\begin{itemize}
    \item Nhận dạng hình ảnh và giọng nói.
    \item Phân loại văn bản, email spam.
    \item Trợ lý áo, chatbot.
    \item Hệ thống gợi ý.
    \item Chẩn đoán y tế, phát hiện gian lận.
\end{itemize}